\clearpage
\fancyhead[L]{\small FADENINTERVALLE.}
\hypertarget{native-450-fadenintervalle}{}
\begin{center}
{\large FADENINTERVALLE.}
\end{center}

\begin{center}
\rule{0.5\linewidth}{0.4pt}\\[2pt]
Astronomische Nachrichten. Band II. Nr. 43. Seite 371\,..\,376. 1823 November.\\[2pt]
\rule{0.5\linewidth}{0.4pt}
\end{center}

\begin{center}
Neue Methode, die gegenseitigen Abstände der Fäden in Meridianfernrohren zu bestimmen.
\end{center}

Die Fadenintervalle in Meridianfernrohren scheint man ausschließlich aus beobachteten Durchgängen des Polarsterns oder anderer polnaher
\SourceLeaf{451}%
Sterne mit bekannter Deklination zu bestimmen. Unter günstigen Bedingungen verbindet dieses Verfahren Bequemlichkeit mit großer Genauigkeit. In Jahreszeiten, in denen ungünstiges Wetter vorherrscht, muss man oft lange warten und manchen vergeblichen Versuch machen, bis man auch nur einen vollständigen Durchgang beobachten kann. Wenn die Gestirne, wie so oft, mit starkem Zittern statt ganz ruhig durch das Gesichtsfeld ziehen, ist die erreichbare Genauigkeit des Verfahrens auch viel geringer. Das alles wäre weniger wichtig, wenn man die endgültige Bestimmung der Intervalle immer auf eine Zeit verschieben könnte, in der die Bedingungen für dieses Verfahren besonders günstig sind. Man möchte aber einerseits die Reduktionen, für die man die Fadenintervalle kennen muss, früh durchführen können. Andererseits kann ein Aufschub bedenklich werden, weil besonders ein Netz aus feinen Spinnenfäden so leicht beschädigt werden kann. Ich selbst musste das Netz mitunter in kurzen Abständen wiederholt erneuern, weil ein Faden erschlafft oder zufällig gerissen war, noch bevor auch nur ein vollständiger Durchgang des Polarsterns bei günstigen Luftverhältnissen beobachtet werden konnte. Auch wenn man die neuen Fäden in dieselben vorgeritzten Kerben legt, kann man keineswegs auf gleiche Abstände zählen, wie mich die Erfahrung oft gelehrt hat. Aus diesen Gründen suchte ich nach einem neuen Verfahren zur Bestimmung der Fadenintervalle, das neben dem üblichen und bei ungünstigen Bedingungen an seiner Stelle mit Vorteil verwendet werden könnte. Das Verfahren, das ich vor Kurzem erprobt habe, scheint seinen Zweck zu erfüllen.

Alle Lichtstrahlen, die parallel zueinander auf das Objektiv treffen, vereinigen sich in einem Punkt der Ebene, in der das Fadennetz liegen soll. Umgekehrt werden alle Strahlen, die in entgegengesetzter Richtung von einem Punkt dieser Ebene ausgehen und auf das Objektiv treffen, nach dem Durchgang durch das Objektiv parallel. Die Strahlen, die von verschiedenen Punkten ausgehen, bilden dagegen nach dem Durchgang durch das Objektiv miteinander genau die Winkel, die dem Abstand dieser Punkte entsprechen, wenn man ihn wie bei der Benutzung des Instruments als Winkel angibt. Ist also die Okularseite des Fernrohrs gegen den Himmel oder eine andere helle Fläche gerichtet, würde ein weitsichtiges Auge das Fadennetz durch das Objektiv deutlich und mit dem richtigen Abstand im eben beschriebenen Sinn sehen, wenn es für so feine Gegenstände empfindlich genug wäre. Was dem bloßen Auge unmöglich ist, gelingt mit einem zweiten Fernrohr, dessen Okular so eingestellt ist, wie es das Auge für sehr weit entfernte Gegenstände braucht. Tatsächlich sieht man so die im Brennpunkt eines \textsc{Frauenhofer}schen Objektivs eingezogenen Fäden mit größter Deutlichkeit und Klarheit.
Auf diese Weise lassen sich auch die Fadenabstände sehr genau messen. Dazu benutzte ich denselben 12-zölligen Theodoliten von \textsc{Ertel}, mit dem ich bei der Gradmessung die Winkel beobachtet habe. Sein eigenes Netz besteht aus einem waagerechten und zwei parallelen Fäden, die eine halbe Minute voneinander entfernt sind. In einem neuen Okular, das \textsc{Repsold} dafür angefertigt hat, sind diese Fäden erstaunlich fein. Man erhält eine sehr saubere Messung, wenn man die durch beide Objektive gesehenen Fäden des Meridianinstruments genau zwischen die beiden parallelen Fäden bringt.

Bei diesem Verfahren muss man das Okular des Meridianfernrohrs selbstverständlich ganz abnehmen oder zumindest eines mit schwächerer Vergrößerung einsetzen, durch das die Fäden, deren Abstände man messen will, gleichzeitig sichtbar sind.

Der Abstand des Theodoliten vom Objektiv spielt zwar keine Rolle, doch wird man ihn lieber nicht zu groß wählen, damit man nicht einige Strahlen zu nah am Rand des Objektivs oder gar nur bei halber Beleuchtung empfängt.

\Needspace{3\baselineskip}
\SourceLeaf{452}%
Der Theodolit gibt, sofern er richtig nivelliert ist, die auf den Horizont projizierten Winkel statt der tatsächlichen Winkel an. Deshalb muss man die unmittelbar gemessenen Werte noch auf die wahren Winkel reduzieren. Bei so kleinen Abständen wie an meinen Meridianinstrumenten genügt dafür die Multiplikation mit dem Kosinus des Neigungswinkels. Bei einem Meridiankreis gibt dieser selbst den Neigungswinkel bereits mit größter Genauigkeit an. Bei einem einfachen Mittagsfernrohr, das die Neigung bestenfalls auf eine Minute genau angibt, kann man sie ebenso gut am kleinen Höhenkreis des Theodoliten ablesen. Dann sollte man eine zu große Neigung vermeiden.

Die folgenden Messungen sind nur als erster Versuch anzusehen. Dabei wurden noch keineswegs alle Vorsichtsmaßnahmen getroffen, die sich empfehlen. Der Theodolit stand mit seinem hölzernen Stativ lediglich auf dem Zwischenboden aus Brettern zwischen den Pfeilern der Meridianinstrumente und war vom Stuhl des Beobachters nicht unabhängig. Die Stellung des Stuhls war bei der beträchtlichen Neigung etwas beengt, und zum Ablesen war das Licht nicht besonders günstig. Wenn man diese Schwierigkeiten vermeidet, hat man umso mehr Grund, eine bessere Übereinstimmung zu erwarten. Jede Zahl ist das Ergebnis von 10 Wiederholungen.

\Needspace{23\baselineskip}
\emph{Abstände des mittleren Fadens im Fernrohr des \textup{\textsc{Reichenbach}}schen Meridiankreises von den 6 übrigen.}

\begin{center}
1823 Okt.\ 23. \quad Neigung $=23^\circ\,58'\,52''$

\begingroup\renewcommand{\arraystretch}{1.1}
\begin{tabular}{r@{\;}c@{}|r|r}
 & & Horizontal & Wahrer Abstand \\
\cline{3-4}
\smash{\raisebox{-0.6em}{I}} & \smash{\raisebox{-0.6em}{$\bigg\{$}} & $11'\,25''050$ & $10'\,25''93$ \\
 & & $11\quad23.000$ & $10\quad24.86$ \\
\smash{\raisebox{-0.6em}{II}} & \smash{\raisebox{-0.6em}{$\bigg\{$}} & $7\quad45.900$ & $7\quad5.69$ \\
 & & $7\quad46.275$ & $7\quad6.03$ \\
\smash{\raisebox{-0.6em}{III}} & \smash{\raisebox{-0.6em}{$\bigg\{$}} & $3\quad58.225$ & $3\quad37.66$ \\
 & & $3\quad57.750$ & $3\quad37.23$ \\
\smash{\raisebox{-0.6em}{V}} & \smash{\raisebox{-0.6em}{$\bigg\{$}} & $3\quad58.850$ & $3\quad38.24$ \\
 & & $3\quad59.375$ & $3\quad38.71$ \\
\smash{\raisebox{-0.6em}{VI}} & \smash{\raisebox{-0.6em}{$\bigg\{$}} & $7\quad53.625$ & $7\quad12.74$ \\
 & & $7\quad53.875$ & $7\quad12.98$ \\
\smash{\raisebox{-0.6em}{VII}} & \smash{\raisebox{-0.6em}{$\bigg\{$}} & $11\quad17.600$ & $10\quad19.11$ \\
 & & $11\quad17.775$ & $10\quad19.28$ \\
\end{tabular}
\endgroup

\medskip
Nov.\ 6. \quad Neigung $=23^\circ\,57'\,30''$

\begingroup\renewcommand{\arraystretch}{1.1}
\begin{tabular}{r@{\;}c@{}|r|r}
\smash{\raisebox{-0.6em}{I}} & \smash{\raisebox{-0.6em}{$\bigg\{$}} & $11'\,24''800$ & $10'\,25''29$ \\
 & & $11\quad25.575$ & $10\quad26.51$ \\
\end{tabular}
\endgroup
\end{center}

Die Fadenintervalle wurden auch aus drei beobachteten Durchgängen des Polarsterns berechnet, von denen sich jedoch nur einer vollständig beobachten ließ. Die Luftverhältnisse waren dabei weder besonders günstig noch besonders ungünstig. Die ersten beiden waren untere Durchgänge bei Tag, der dritte ein oberer Durchgang bei Nacht.

\begin{center}
\begingroup\small\setlength{\tabcolsep}{4pt}
\begin{tabular}{l|r|r|r|r|r|r}
 & I & II & III & V & VI & VII \\
Okt.\ 23 & $10'\,26''91$ & .\,.\,.\,. & $3'\,38''54$ & $3'\,38''12$ & $7'\,11''27$ & $10'\,19''31$ \\
Nov.\ 5 & $10\quad26.90$ & $7'\,7''01$ & $3\quad38.33$ & $3\quad39.39$ & $7\quad10.41$ & $10\quad19.50$ \\
\phantom{Nov.\ }6 & $10\quad24.31$ & .\,.\,.\,. & $3\quad35.12$ & $3\quad40.22$ & .\,.\,.\,. & .\,.\,.\,. \\
\end{tabular}
\endgroup
\end{center}

Wendet man die in meiner Theoria Combinationis Observationum entwickelten Grundsätze auf die Bestimmungen mit dem Theodoliten an, so ergibt sich für ein Fadenintervall aus 10 Wiederholungen ein mittlerer Fehler $= 0''67$. Bei Anwendung auf die Durchgänge des Polarsterns ergibt sich aus einem Durchgang dagegen ein mittlerer Fehler $= 1''38$. Die erste Bestimmung hätte demnach eine doppelt so große Genauigkeit oder das vierfache Gewicht. Allerdings ist die Zahl der Messungen
\SourceLeaf{453}%
auf beiden Seiten zu gering, um daraus den mittleren Fehler mit großer Zuverlässigkeit abzuleiten. Beim Polarstern gilt er ohnehin nur für ähnliche Luftverhältnisse, da er bei außergewöhnlich günstigen Luftverhältnissen bedeutend kleiner ausfallen würde. Im Allgemeinen bestätigt sich dadurch aber zumindest die Brauchbarkeit des neuen Verfahrens hinreichend.

Eine andere, nicht weniger interessante Anwendung der obigen Grundsätze ist die Bestimmung der Vergrößerung von Fernrohren, die ich hier aber nur kurz andeute. Ist das Okular für ein weitsichtiges Auge eingestellt, erscheinen die Gegenstände beim Blick von hinten durch das Fernrohr um denselben Faktor verkleinert, um den es sie beim üblichen Gebrauch vergrößert. Sie sind dabei so deutlich zu sehen, dass man sie ohne Weiteres durch ein zweites Fernrohr wieder vergrößern kann. Man kann also den Winkel zwischen zwei auf diese Weise betrachteten Objekten mit dem Theodoliten durch Wiederholung messen und ihn mit dem Winkel zwischen den Richtungen von der Mitte des Okulars zu den Objekten selbst vergleichen. Die Objekte brauchen gar nicht sehr weit entfernt zu sein. Bei etwas stärkerer Vergrößerung genügen unter Umständen wenige Fuß. Für größere Genauigkeit muss man bei einigermaßen großen Winkeln am Okular gebührend berücksichtigen, dass die Vergrößerungen genauer genommen zu den Tangenten der Neigungswinkel gegen die optische Achse proportional sind und nicht zu den Winkeln selbst. Am einfachsten wählt man also die Objekte so, dass die Richtung der optischen Achse sehr nahe in der Mitte zwischen den Richtungen von der Mitte des Okulars zu diesen Objekten liegt. Dann vergleicht man die Tangente des halben Winkels am Okular mit der Tangente des halben Winkels, der sich nach gehöriger Reduktion der Theodolitenbeobachtung auf den wahren Wert ergibt. Statt des Theodoliten kann man auch das Heliometer verwenden, das den wahren Wert unmittelbar angibt.
Dieses Verfahren ermöglicht eine ungleich größere Genauigkeit als die anderen, die man sonst für denselben Zweck angewandt hat.
